\chapter{Lesson 11: Term Review}

The last week of the term was spent consolidating grammar before the final exam (oral
interview, listening, character recognition and a short composition). These are my
summary notes.

\section{Asking questions}

\begin{center}
\begin{tikzpicture}[c/.style={draw, rounded corners, align=center, font=\small, fill=blue!6, text width=3.3cm}, h/.style={draw, very thick, rounded corners, fill=yellow!25, font=\bfseries}]
  \node[h] (Q) at (0,0) {Questions};
  \node[c] (ma) at (-5,2.1) {\zh{吗}\\ statement + \zh{吗}\\ \zh{你是学生吗？}};
  \node[c] (an) at (0,2.6) {A-not-A\\ \zh{你忙不忙？}\\ no \zh{吗}, no \zh{也/都}};
  \node[c] (qw) at (5,2.1) {question word in place\\ \zh{什么, 谁, 哪, 哪儿, 几, 多少, 怎么, 为什么}};
  \node[c] (hs) at (-5,-2.1) {choice: A \zh{还是} B\\ \zh{茶还是咖啡？}};
  \node[c] (ne) at (0,-2.6) {follow-up: X \zh{呢}？\\ \zh{我很好，你呢？}};
  \node[c] (ba) at (5,-2.1) {confirming: \zh{吧}\\ \zh{你是小王吧？}\\ also \zh{好吗/怎么样/是不是}};
  \foreach \n in {ma,an,qw,hs,ne,ba} \draw[thick, gray] (Q) -- (\n);
\end{tikzpicture}
\end{center}

\begin{important}[Never double up]
Only one question device per sentence: no \zh{吗} with A-not-A, with a question word or
with \zh{还是}.
\end{important}

\section{Word order}

\begin{enumerate}
  \item Subject -- Time -- Adverb -- Prepositional phrase (\zh{在/跟/给}) -- Verb --
        Object.
  \item Time-when goes before the verb; time clauses with \zh{以前/以后/的时候} come first.
  \item Adverbs (\zh{也, 都, 常常, 就, 才, 不, 没}) go right before the verb or the
        prepositional phrase. Order when several appear: \zh{也都不}.
  \item Degree complements go after the verb: V \zh{得} adj.
\end{enumerate}

\begin{example}[Sorting jumbled words]
Put in order: \zh{在 / 我们 / 星期六 / 图书馆 / 一起 / 中文 / 复习 / 常常 / 跟同学}.

Answer: \zh{我们星期六常常跟同学一起在图书馆复习中文。}
\end{example}

\section{The particle \zh{的} so far}

\begin{center}
\begin{tabular}{llp{6cm}}
\toprule
Use & Example & Notes \\
\midrule
possession & \zh{老师的书} & dropped with pronoun + kin: \zh{我妈妈} \\
description & \zh{中国的电影}, \zh{很好的朋友} & adjective + \zh{的} if more than one syllable or with \zh{很} \\
noun omitted & \zh{这本书是我的。} & ``mine'' \\
clauses + \zh{的时候} & \zh{吃饭的时候} & ``when'' \\
\zh{是\ldots{}的} (preview) & \zh{你是怎么来的？} & emphasises time/place/manner of a past event \\
\bottomrule
\end{tabular}
\end{center}

One-syllable adjectives usually attach without \zh{的}: \zh{好朋友}, \zh{新电脑},
\zh{大学校}.

\section{Negation}

\begin{center}
\begin{tabular}{lp{6cm}l}
\toprule
Negator & Used with & Example \\
\midrule
\zh{不} & present/future, habits, adjectives, \zh{是}, modals & \zh{我不去}, \zh{不忙}, \zh{不是}, \zh{不会} \\
\zh{没(有)} & \zh{有}; completed actions (drop \zh{了}) & \zh{没有弟弟}, \zh{我没去} \\
\zh{别} & commands & \zh{别等我} \\
\zh{不用} & ``don't have to'' (negates \zh{得}) & \zh{你不用来} \\
\bottomrule
\end{tabular}
\end{center}

\section{The particle \zh{了} so far}

\begin{itemize}
  \item After the verb: the action took place / was completed.
        \zh{我买了两本书。} Needs a quantified object or a following clause.
  \item At the end of a sentence (preview of CHINA 102R): a change of situation or
        ``new information''. \zh{我不喝咖啡了。} I don't drink coffee any more.
  \item \zh{太\ldots{}了} is a fixed pattern, not aspect.
  \item Negative: \zh{没} + verb, no \zh{了}.
  \item Not used with \zh{才}, with habitual past actions, or with verbs like \zh{是},
        \zh{在}, \zh{喜欢}, \zh{觉得} in the past.
\end{itemize}

\section{Measure words}

\begin{center}
\begin{tabular}{llll}
\toprule
MW & Pinyin & For & Example \\
\midrule
\zh{个} & gè & general, people & \zh{一个朋友} \\
\zh{位} & wèi & people (polite) & \zh{那位先生} \\
\zh{口} & kǒu & family size & \zh{五口人} \\
\zh{张} & zhāng & flat things & \zh{一张照片} \\
\zh{本} & běn & books & \zh{三本书} \\
\zh{杯} & bēi & cup of & \zh{一杯茶} \\
\zh{瓶} & píng & bottle of & \zh{两瓶可乐} \\
\zh{门} & mén & courses & \zh{五门课} \\
\zh{节} & jié & class periods & \zh{两节课} \\
\zh{只} & zhī & animals & \zh{一只猫} \\
\bottomrule
\end{tabular}
\end{center}

\section{Common errors from the midterm}

\begin{center}
\begin{tabular}{lll}
\toprule
Wrong & Right & Rule \\
\midrule
*\zh{我是很忙。} & \zh{我很忙。} & no \zh{是} before adjectives \\
*\zh{我工作在银行。} & \zh{我在银行工作。} & place before verb \\
*\zh{也我是学生。} & \zh{我也是学生。} & \zh{也} after subject \\
*\zh{我有二个哥哥。} & \zh{我有两个哥哥。} & \zh{两} before MW \\
*\zh{我不有车。} & \zh{我没有车。} & \zh{有} takes \zh{没} \\
*\zh{他说中文得很好。} & \zh{他说中文说得很好。} & repeat the verb \\
*\zh{你忙不忙吗？} & \zh{你忙不忙？} & no double question \\
*\zh{我没去了。} & \zh{我没去。} & no \zh{了} after \zh{没} \\
*\zh{这个字一点儿难。} & \zh{这个字有一点儿难。} & \zh{有一点儿} + adj \\
\bottomrule
\end{tabular}
\end{center}

\section{Self-introduction for the oral exam}

\begin{example}[Model answer]
\zh{大家好！我姓陈，叫陈安。我是加拿大人，今年十九岁。我是滑铁卢大学一年级的学生，我的专业是电脑。我家有四口人：爸爸、妈妈、姐姐和我。我爸爸是工程师，妈妈是老师。我姐姐在多伦多工作。我很喜欢学中文，因为我觉得中文很有意思。我每天早上八点起床，上午上课，下午常常跟朋友一起去图书馆复习。周末我喜欢打球、看电影。谢谢！}

(\zh{多伦多} Duōlúnduō = Toronto.)
\end{example}

\section*{Review}
\begin{itemize}
  \item Six question types; never combine two.
  \item Word order: S -- T -- Adv -- PP -- V -- O (-- \zh{得} complement).
  \item \zh{的} for possession and description; one-syllable adjectives attach directly.
  \item \zh{不} vs \zh{没} vs \zh{别} vs \zh{不用}.
  \item \zh{了} after the verb = completed action; sentence-final \zh{了} = change (next
        term).
  \item Know the measure words and the common error list.
\end{itemize}
